% ch07 -- Informationseinheit und Skala
\label{ch:informationseinheit_skala}
% Quelldokument: 257 -- Ableitung hier vollständig geführt

\epigraph{$k_BT\ln 2$ ist kein universelles Gesetz.
Es ist der Spezialfall einer bestimmten Skala und Temperatur.}{Dok.\,257}

\section{Die Frage}

Woher kommt die Zahl $k_BT\ln 2$ in Landauers Formel? Warum
genau diese Kombination aus Temperatur, Boltzmann-Konstante und dem Logarithmus
von 2? Und gilt sie überall --- oder nur unter bestimmten Bedingungen?

Die FFGFT gibt eine präzise Antwort: Sie gilt genau dann, wenn die
geometrische Energie eines Windungsquants auf der betrachteten Skala
mit der thermischen Energie des Bades übereinstimmt. Das ist eine
Bedingung, keine Universalkonstante.

Die vollständige Ableitung folgt in drei Ebenen. \status{B Dok.\,257}

\section{Ebene 1: Das Windungsquant --- topologisch, universell}

\begin{keyresult}[Minimale Informationseinheit --- \status{SETZUNG}]
Die minimale Informationseinheit in der FFGFT ist das \textbf{Windungsquant}
$\Delta w = 1$ auf dem fundamentalen Torus $T^4$.

Die Homotopiegruppe des Torus ist $\pi_1(T^4) \cong \mathbb{Z}^4$ ---
sie kennt nur ganzzahlige Windungszahlen. Ein halbiertes Windungsquant
existiert nicht. Diese Diskretheit ist eine topologische Tatsache,
kein Postulat. \status{Q}
\end{keyresult}

\textit{Zur externen Absicherung:} $\pi_1(T^n) \cong \mathbb{Z}^n$ ist ein
Standardresultat der algebraischen Topologie. Es folgt aus dem Satz, dass
der $n$-dimensionale Torus $T^n = S^1 \times \cdots \times S^1$ ($n$-faches
Produkt von Kreisen) ist, und dass $\pi_1(S^1) \cong \mathbb{Z}$ (der
Kreis hat Windungszahl $\in \mathbb{Z}$), und dass die Fundamentalgruppe
eines Produktraums das Produkt der Fundamentalgruppen ist:
\[
  \pi_1(T^4) = \pi_1(S^1)^4 \cong \mathbb{Z}^4.
\]
Lehrbuch-Referenzen: Hatcher, \textit{Algebraic Topology}~\cite{Hatcher2002},
Kap.~1.2 (frei verfügbar: \url{https://pi.math.cornell.edu/~hatcher/AT/AT.pdf});
Bredon, \textit{Topology and Geometry}~\cite{Bredon1993}, Kap.~III.§5.
Das FFGFT-spezifische Argument --- Identifikation von $\Delta w = 1$ mit
1 Bit --- ist die konzeptuelle Setzung; die Ganzzahligkeit selbst ist
externe, etablierte Mathematik. \status{Q}

Die Identifikation $\Delta w = 1 \leftrightarrow$ 1 Bit ist die konzeptuelle
Brücke zwischen Topologie und Informationstheorie. Sie ist
\textbf{skalenunabhängig}: Ob auf der Planck-Skala oder der Zellskala ---
ein Windungsquant ist ein Windungsquant. \status{B}

\begin{laierspur}
Stell dir eine Schnur vor, die einmal um einen Pfosten gewickelt ist.
Du kannst sie zweimal wickeln oder dreimal --- aber nicht anderthalb Mal,
ohne dass sie aufgeht. Das ist $\pi_1(T^4) \cong \mathbb{Z}^4$:
Windungszahlen sind ganzzahlig, erzwungen durch die Topologie des Torus.
Ein Bit entspricht genau einer solchen Windung.
\end{laierspur}

\section{Ebene 2: Die Bit-Energie --- geometrisch, skalenabhängig}

Ein Windungsquant auf einer physikalischen Skala $L$ hat eine Energie.
Diese Energie folgt aus zwei bekannten Prinzipien:

\paragraph{Schritt 1 --- Wirkung des Windungsquants.}
Die Wirkung eines Windungsquants beträgt $S = \Delta w \cdot h = h$.
Das ist eine direkte Konsequenz der Quantisierungsbedingung auf $T^4$.

\paragraph{Schritt 2 --- Impuls aus der Unschärferelation.}
Auf einer räumlichen Skala $L$ gilt $\Delta x \cdot \Delta p \sim \hbar$,
also $p \sim \hbar/L$.

\paragraph{Schritt 3 --- Energie des Windungsquants.}
Mit $E = pc$ (masselose Mode):
\[
  \boxed{E_\text{bit}(L) = \frac{\hbar c}{L}}
\]

Das ist die Energie eines minimalen Informationsquants auf der Skala $L$.
Sie hängt von der Skala ab --- und nur von ihr. \status{B}

\begin{laierspur}
Die Formel $\hbar c / L$ ist aus der Physik vertraut: Das ist die Energie
eines Photons mit Wellenlänge $L$. Ein Windungsquant auf dem Torus verhält
sich energetisch wie ein solches Photon --- seine Energie ist durch die
räumliche Skala bestimmt, auf der es existiert.
\end{laierspur}

\section{Ebene 3: Landauer als Spezialfall --- thermodynamisch, eine Skala}

Landauers Formel $E_L = k_BT\ln 2$ gilt im thermodynamischen Regime.
Das ist genau dann der Fall, wenn die geometrische Bit-Energie mit
der thermischen Energie des Bades übereinstimmt:
\[
  E_\text{bit}(L) = k_BT\ln 2
  \quad\Longleftrightarrow\quad
  \frac{\hbar c}{L} = k_BT\ln 2
\]
Auflösen nach der Skala ergibt die \textbf{Äquivalenztemperatur}:
\[
  T_\text{äquiv}(L) = \frac{\hbar c}{k_B \, L \, \ln 2}
\]
Das ist die Temperatur, bei der Landauers Formel für eine gegebene Skala $L$
gilt --- und \emph{nur} bei dieser Temperatur. \status{B}

\section{Die Konsequenz: 70 Größenordnungen Variation}

Setzt man verschiedene physikalische Skalen ein:

\begin{center}
\begingroup\small
\begin{tabular}{@{}p{2.5cm}p{2.5cm}p{3.0cm}@{}}
\toprule
Skala $L$ & $E_\text{bit}(L)$ & $T_\text{äquiv}$ [K] \\
\midrule
Planck ($10^{-35}$\,m) & $E_P \approx 1{,}96$\,GJ & $2 \times 10^{32}$ \\
Proton ($10^{-15}$\,m) & $\sim 200$\,MeV & $3 \times 10^{12}$ \\
Atom ($10^{-10}$\,m) & $\sim 2$\,keV & $3 \times 10^{7}$ \\
\textbf{Zelle} ($10^{-5}$\,m) & $\sim 0{,}02$\,eV & $\mathbf{\approx 331}$ \\
p-Bit ($22$\,nm) & $9{,}2$\,meV & $\approx 100$ \\
Kosmisch ($10^{26}$\,m) & $\sim 10^{-51}$\,J & $10^{-30}$ \\
\bottomrule
\end{tabular}
\endgroup
\end{center}

Landauers $k_BT\ln 2$ bei Raumtemperatur ($T = 300$\,K) liegt bei
${\approx}2{,}9$\,meV --- das entspricht einer Skala von ${\approx}48\,\mu$m,
nahe der Zellskala. \status{K}

\begin{laierspur}
Die Tabelle zeigt: Die Äquivalenztemperatur variiert über \emph{70 Größenordnungen}.
Das bedeutet: Landauers $k_BT\ln 2$ gilt exakt nur bei einer einzigen Skala
für eine gegebene Temperatur --- nicht überall. Es ist kein universelles Gesetz,
sondern ein Treffer im $E_\text{bit}(L)$-Spektrum.
\end{laierspur}

\section{Die Hierarchie --- zusammengefasst}

\[
  T^4\text{-Geometrie}
  \;\longrightarrow\; \Delta w = 1
  \;\longrightarrow\; E_\text{bit}(L) = \frac{\hbar c}{L}
  \;\longrightarrow\; E_L = k_BT\ln 2 \;\text{(Spezialfall)}
\]

Geometrie ist primärer als Energie. Energie ist eine skalenabhängige
Projektion der $T^4$-Geometrie, nicht ihr Fundament. \status{B}

Matzkes Ansatz setzt $T = T_P$ und bekommt:
\[
  \mbit = \frac{\mP \ln 2}{2\pi} \approx 0{,}11\,\mP
\]
Das ist der Treffer bei der Planck-Skala --- eine Setzung, die man
im FFGFT-Bild auch rechnen kann: $E_\text{bit}(L_P) = E_P$.
Der Unterschied: Matzke deklariert diesen Treffer als universell.
FFGFT zeigt, dass er skalensensitiv ist. \status{K}

\section{Was bewiesen und was offen ist}

\textbf{Bewiesen \status{B}:} Die Hierarchie Topologie $\to$ Geometrie $\to$
Thermodynamik ist intern konsistent. Die Ableitung $E_\text{bit}(L) = \hbar c/L$
aus Windungsquantisierung und Unschärferelation ist mathematisch korrekt.
Landauers Formel ist in dieser Hierarchie ein Spezialfall.

\textbf{Hergeleitet und numerisch überprüft \status{K}:} Die Äquivalenztemperatur
stimmt auf der Zellskala mit biologischer Raumtemperatur überein.
Die p-Bit-Vorhersage $L_8 \approx 22$\,nm ist konsistent.

\textbf{Offen \status{S}:} Gibt es einen direkten experimentellen Test, der
$E_\text{bit}(L) = \hbar c/L$ von Landauers $k_BT\ln 2$ unterscheidet?
Das wäre ein Experiment auf einer Skala, bei der $T_\text{äquiv}$ weit
von der Badtemperatur entfernt ist.
